\section{Latihan Soal Tambahan: Logika Proposisi}

Berikut adalah kumpulan soal latihan untuk memperdalam pemahaman tentang proposisi, operator logika, tabel kebenaran, dan evaluasi ekspresi logika.

\subsection{Soal Konversi Bahasa Alami ke Simbol Logika}

\textbf{Soal 1.}
Gunakan konstanta proposisional berikut:
\begin{center}
\begin{tabular}{lcl}
  $p$ & = & Bowo kaya raya \\
  $q$ & = & Bowo hidup bahagia
\end{tabular}
\end{center}
Ubah pernyataan-pernyataan berikut menjadi bentuk logika:
\begin{enumerate}
  \item Bowo tidak kaya.
  \item Bowo kaya raya dan hidup bahagia.
  \item Bowo kaya raya atau tidak hidup bahagia.
  \item Jika Bowo kaya raya, maka ia hidup bahagia.
  \item Bowo hidup bahagia jika dan hanya jika ia kaya raya.
\end{enumerate}

\textbf{Soal 2.}
Misalkan $p$, $q$, dan $r$ adalah variabel proposisional:
\begin{center}
\begin{tabular}{lcl}
  $p$ & = & Anda sakit flu \\
  $q$ & = & Anda ujian \\
  $r$ & = & Anda lulus
\end{tabular}
\end{center}
Ubahlah ekspresi berikut menjadi pernyataan dalam bahasa Indonesia:
\begin{enumerate}
  \item $p \rightarrow \neg q$
  \item $q \rightarrow \neg r$
  \item $\neg q \rightarrow r$
  \item $(p \wedge q) \rightarrow r$
  \item $(p \rightarrow \neg r) \vee (q \rightarrow \neg r)$
  \item $(p \wedge q) \vee (\neg q \wedge r)$
\end{enumerate}

\textbf{Soal 3.}
Ubahlah pernyataan-pernyataan berikut menjadi bentuk logika (definisikan variabel proposisional yang digunakan):
\begin{enumerate}
  \item Jika Bowo berada di Malioboro, maka Dewi juga ada di Malioboro.
  \item Pintu rumah Dewi berwarna merah atau coklat.
  \item Berita itu tidak menyenangkan.
  \item Bowo akan datang jika ia mempunyai kesempatan.
  \item Jika Dewi rajin kuliah, maka ia pasti pandai.
\end{enumerate}

\subsection{Soal Tabel Kebenaran}

\textbf{Soal 4.}
Dengan menggunakan tabel kebenaran, jawab pertanyaan-pertanyaan berikut:
\begin{enumerate}
  \item Apakah nilai kebenaran dari $p \wedge p$?
  \item Apakah nilai kebenaran dari $p \vee p$?
  \item Apakah nilai kebenaran dari $(p \wedge \neg p)$ dan $(p \vee \neg p)$?
  \item Apakah $(p \rightarrow q)$ mempunyai nilai kebenaran yang sama dengan $(q \rightarrow p)$?
  \item Apakah $(p \rightarrow q) \rightarrow r$ mempunyai nilai kebenaran yang sama dengan $p \rightarrow (q \rightarrow r)$?
\end{enumerate}

\textbf{Soal 5.}
Buatlah tabel kebenaran untuk ekspresi-ekspresi logika berikut:
\begin{enumerate}
  \item $\neg(\neg p \wedge \neg q)$
  \item $p \wedge (p \vee q)$
  \item $((\neg p \wedge (\neg q \wedge r)) \vee (q \wedge r)) \vee (p \wedge r)$
  \item $(p \wedge q) \vee (((\neg p \wedge q) \rightarrow p) \wedge \neg q)$
  \item $(p \rightarrow q) \leftrightarrow (\neg q \rightarrow \neg p)$
  \item $p \wedge ((r \vee q) \leftrightarrow \neg r)$
  \item $\neg((p \wedge q) \rightarrow \neg r) \vee p$
\end{enumerate}

\subsection{Soal Konversi ke Proposisi Majemuk}

\textbf{Soal 6.}
Ubahlah pernyataan-pernyataan berikut menjadi ekspresi logika berupa proposisi majemuk:
\begin{enumerate}
  \item Jika tikus itu waspada dan bergerak cepat, maka kucing atau anjing itu tidak mampu menangkapnya.
  \item Jika saya tidak keliru, Dewi sudah diwisuda dan pacarnya atau orangtuanya berada disampingnya.
  \item Bowo membeli saham dan membeli properti untuk investasinya, atau dia dapat menanamkan uang di deposito bank dan menerima bunga uang.
\end{enumerate}

\subsection{Soal Hirarki Operator}

\textbf{Soal 7.}
Masukkan tanda kurung biasa ke dalam ekspresi logika berikut ini sehingga tidak terjadi ambiguitas:
\begin{enumerate}
  \item $p \wedge q \wedge r \rightarrow s$
  \item $p \vee q \vee r \leftrightarrow \neg s$
  \item $\neg p \wedge q \rightarrow \neg r \vee s$
  \item $p \rightarrow q \leftrightarrow \neg r \rightarrow \neg s$
  \item $p \vee q \wedge r \rightarrow p \wedge q \vee r$
\end{enumerate}

\subsection{Soal Evaluasi Ekspresi Logika}

\textbf{Soal 8.}
Jika $\tau(p) = B$ dan $\tau(q) = B$, sedangkan $\tau(r) = S$ dan $\tau(s) = S$, carilah nilai kebenaran dari ekspresi-ekspresi logika berikut:
\begin{enumerate}
  \item $p \wedge (q \vee r)$
  \item $(p \vee q) \wedge r$
  \item $((p \vee q) \wedge r) \vee \neg((p \vee q) \wedge (q \vee s))$
  \item $(\neg(p \wedge q) \vee \neg r) \vee (((\neg p \wedge q) \vee \neg s) \wedge r)$
  \item $(p \leftrightarrow r) \wedge (\neg q \rightarrow s)$
\end{enumerate}

\textbf{Soal 9.}
Jika $\tau(p) = S$ dan $\tau(q) = S$, sedangkan $\tau(r) = B$ dan $\tau(s) = B$, carilah nilai kebenaran dari ekspresi-ekspresi logika berikut:
\begin{enumerate}
  \item $\neg(p \leftrightarrow q) \wedge (\neg r \rightarrow s)$
  \item $(\neg p \wedge (q \vee \neg r) \vee (q \leftrightarrow \neg p)) \leftrightarrow (s \wedge r)$
  \item $(p \vee (q \rightarrow (r \wedge \neg p))) \leftrightarrow \neg(q \vee \neg s)$
\end{enumerate}
